\BourbakiExerciseGroup

\begin{enumerate}
\def\labelenumi{\arabic{enumi})}
\tightlist
\item
  设 K 为一个域，且 \((\mathrm{E}_{i})_{i\in I}\) 为 K 的任意一族扩张。证明存在 K 的一个扩张 E，并且对每个 \(i\in I\) ，存在一个 K-同构 \(u_{i}\) （从 \(E_{i}\) 到 E 中的），使得 E 由诸 \(u_{i}(\mathrm{E}_{i})\) 的并生成（如命题 4（V，第 22 页）中那样论证）。
\end{enumerate}

* 2) 设 \(\mathbf{K}\) 是特征为 0 的代数闭域， \(\mathrm{F} = \mathrm{K}((\mathbf{X}^{1 / n!}))\) 为 \(\mathbf{X}^{1 / n!}\) 的形式幂级数域，而 \(\mathrm{E} = \mathrm{K}((\mathbf{X}))\) ，其中 \(\mathbf{X} = (\mathbf{X}^{1 / n!})^{n!}\) ，是 \(\mathbf{F}\) 中由 \(\mathbf{X}\) 的形式幂级数组成的子域。对于每个形式幂级数 \(f \in \mathbf{K}((\mathbf{X}))\) ，我们用 \(\omega(f)\) 表示它的阶。

\begin{enumerate}
\def\labelenumi{\alph{enumi})}
\item
  设 \(P(Y) = a_{0}Y^{n} + a_{1}Y^{n-1} + \cdots + a_{n}\) 是 E{[}Y{]} 中次数为 n 的多项式，使得 \(a_{n} \neq 0\) 。证明存在整数的一个严格递增序列 \((i_{k})_{0 \leqslant k \leqslant r}\) （在区间 \([0, n]\) 中），使得： \(1^{\circ} i_{0} = 0\) , \(i_{r} = n\) ; \(2^{\circ} \omega(a_{i_{k}})\) 对 \(0 \leqslant k \leqslant r\) 是有限的; \(3^{\circ}\) 对于每个满足 \(0 \leqslant j \leqslant n\) 、异于诸 \(i_{k}\) 且使 \(\omega(a_{j})\) 为有限的指标 j，点 \((j, \omega(a_{j})) \in \mathbf{R}^{2}\) 位于经过点 \((i_{k}, \omega(a_{i_{k}}))\) , \((i_{k-1}, \omega(a_{i_{k-1}}))\) 的直线上方，且当 \(j < i_{k-1}\) 或 \(j > i_{k} (1 \leqslant k \leqslant r)\) 时严格位于其上方。连接点 \((i_{k-1}, \omega(a_{i_{k-1}}))\) 与 \((i_{k}, \omega(a_{i_{k}}))\) （对 \(1 \leqslant k \leqslant r\) ）的诸线段之并称为 P 的牛顿多边形，前述线段称为其边，而点 \((i_{k}, \omega(a_{i_{k}}))\) 为其顶点。
\item
  令 \(\rho_{k} = i_{k} - i_{k - 1}\) , \(\sigma_{k} = (\omega (a_{i_{k}}) - \omega (a_{i_{k - 1}})) / \rho_{k}\) ；证明 \(P\) 恰有 \(\rho_{k}\) 个零点（按正确的重数计算）属于域 \(K((X^{1 / (\rho_k!)}))\) 且其阶等于 \(\sigma_{k}\) （用递归方法确定满足方程 \(P(z) = 0\) 的这种类型的形式幂级数的系数）.
\item
  得出结论：诸域 \(\mathbf{K}((\mathbf{X}^{1 / n}))\) （对所有整数 \(n > 0\) ）的并集是域 \(\mathrm{E} = \mathbf{K}((\mathbf{X}))\) 的一个代数闭包（“皮瑟定理”）。
\item
  如果 \(\mathbf{K}'\) 是特征为 2 的域，并且 \(\mathbf{E}' = \mathbf{K}'((\mathbf{X}))\) ，则多项式 \(\mathbf{Y}^2 + \mathbf{X}\mathbf{Y} + \mathbf{X}\) （ \(\mathbf{E}'[\mathbf{Y}]\) 中的）在诸域 \(\mathbf{K}'((\mathbf{X}^{1/n!}))\) 的并集中没有根。 *
\end{enumerate}

